\section{A genuine coefficient-two obstruction for the quadratic seven factor}
\label{sec:quadratic-seven-coefficient-two-obstruction}

\paragraph{The original remaining label.}
Put $K=\mathbb Q(\omega)=\mathbb Q(\sqrt{-3})$.
Let $\xi=\chi_{-7}$ be the quadratic Dirichlet character modulo seven.
It is the order-two character missing from the coefficient-two
families of AC-G182 and AC-G184.
Let $E/\mathbb Q_2$ be any finite extension with its original
ring of integers $O$.  For $e\geq1$ put
\[
       a_e=\xi c_2^{-e},\qquad D_e=(E/O)(a_e),
       \qquad b_2=[E:\mathbb Q_2].
 \tag{G186.1}
\]
The coefficient ring has trivial Galois action; only the displayed
scalar acts.  In the more general proposed family, this is the
permitted specialization $\delta=1$ and $\kappa=1$.
Write $\mathcal S_2(K,D_e)$ for full local zero at the unique
place above two and zero on inertia at every other finite place.

\begin{theorem}[AC-G186: an original-integral odd-weight obstruction]
\label{AC-G186}
For every odd positive $e$ there is an explicit $O$-linear injection
\[
                 O/2O\lhook\joinrel\longrightarrow
                           \mathcal S_2(K,D_e).
 \tag{G186.2}
\]
The same injection lands in the actual group
$\mathcal S_{2,\xi}(e)$ of (FC.2), on its unchanged original
completion and with all added host inertia conditions.
Consequently that actual strict group is nonzero.  In the finite
setting of (FC.2),
\[
 \#\mathcal S_{2,\xi}(e)\geq2^{b_2},\qquad
 \frac{\log\#\mathcal S_{2,\xi}(e)}{b_2}\geq\log2.
 \tag{G186.3}
\]
Thus the remaining quadratic-seven family cannot have strict zero
for every two-power label and every positive weight.
This theorem gives a lower bound only.  It gives no upper bound,
no positive derived-index lower bound and no obstruction to a
separately proved bounded-cost window estimate.
\end{theorem}

\begin{proof}
\proofstep{1. The exact quadratic character and its local value at two.}
For a primitive seventh root $\zeta$, set
\[
 A=\zeta+\zeta^2+\zeta^4,\qquad
 B=\zeta^3+\zeta^5+\zeta^6.
\]
The relations $A+B=-1$ and $A^2=A+2B$ give
$A^2+A+2=0$ and $(2A+1)^2=-7$.
An automorphism $\zeta\mapsto\zeta^r$ preserves $A,B$ when
$r$ is a square modulo seven and exchanges them otherwise.
Thus $2A+1$ transforms through the quadratic character modulo
seven, and its field is exactly $\mathbb Q(\sqrt{-7})$.
This proves the stated character-field convention directly.
Its restriction to $G_K$ is nontrivial: the distinct quadratic
fields $\mathbb Q(\sqrt{-7})$ and $K$ intersect in $\mathbb Q$.

The polynomial $f(X)=X^2+X+2$ has a root in $\mathbb Z_2$.
For completeness, start with $x_1=0$ modulo two.
If $f(x_n)=0$ modulo $2^n$, replace $x_n$ by
$x_n+\epsilon2^n$, with $\epsilon\in\{0,1\}$ chosen so that
\[
 f(x_n+\epsilon2^n)
       \equiv f(x_n)+\epsilon2^n(2x_n+1)\pmod{2^{n+1}}
\]
vanishes.  The derivative is odd, so this choice exists and
is unique.  The compatible sequence has a limit $x\in\mathbb Z_2$
with $f(x)=0$.  Hence $(2x+1)^2=-7$ in $\mathbb Q_2$.
It follows that
\[
                     \xi|_{G_{K_2}}=1.
 \tag{G186.4}
\]
In particular a cochain depending only on the quadratic
seven character may vanish on the entire local group at two.

\proofstep{2. A local cyclotomic scalar with exact valuation one.}
The field $K_2=\mathbb Q_2(\sqrt{-3})$ does not contain $i$.
One can check this without a global disjointness assumption.
If $i=r+s\sqrt{-3}$ with $r,s\in\mathbb Q_2$, squaring gives
$2rs=0$.  If $s=0$ then $r^2=-1$, impossible since an odd
two-adic unit square is one modulo eight.
If $r=0$ then $s^2=1/3$, also impossible since $1/3\equiv3$
modulo eight.  The same unit-square criterion shows that
$-3$ is not a square in $\mathbb Q_2$, so the quadratic basis
used here is valid.
The quadratic extension $K_2(i)/K_2$ therefore has a nontrivial
automorphism, which lifts to $g_0\in G_{K_2}$.
By the definition of the cyclotomic character,
\[
              c_2(g_0)\equiv-1\pmod4.
 \tag{G186.5}
\]
For odd $e$, equations (G186.4)--(G186.5) give
\[
                  a_e(g_0)-1=2u,
                        \qquad u\in\mathbb Z_2^\times.
 \tag{G186.6}
\]
The factor $u$ remains a unit in every original $O$.

\proofstep{3. The explicit cocycles on the original module.}
For $t\in D_e[2]$, define
\[
 z_t(g)=\epsilon(g)t,
          \qquad \epsilon(g)=\frac{1-\xi(g)}2\in\{0,1\}.
 \tag{G186.7}
\]
Every value of $a_e$ is a rational two-adic odd unit, so it
acts trivially on $D_e[2]$.
The sign character satisfies
$\epsilon(gh)\equiv\epsilon(g)+\epsilon(h)\pmod2$.
Since $2t=0$, this proves the cocycle identity
$z_t(gh)=z_t(g)+a_e(g)z_t(h)$.
The cocycle is continuous and has finite image.
The assignment $t\mapsto[z_t]$ is $O$-linear.

Equation (G186.4) makes $z_t$ identically zero on $G_{K_2}$.
At a place over seven, the two-adic cyclotomic character is
unramified, so $a_e|_{I_v}=\xi|_{I_v}$.
Multiplication by two on $E/O$ is surjective; choose
$d_t\in E/O$ with $-2d_t=t$.
Then for every $g\in I_v$,
\[
                       z_t(g)=(a_e(g)-1)d_t.
 \tag{G186.8}
\]
Indeed both sides are zero for $\xi(g)=1$ and equal $t$
for $\xi(g)=-1$.
At every other finite place away two, both the quadratic
seven character and the cyclotomic coefficient are unramified,
and $z_t$ is identically zero on inertia.
Thus $[z_t]\in\mathcal S_2(K,D_e)$.
The choice of $d_t$ in this step is only a local coboundary
witness; it is not assumed to be a global witness.

\proofstep{4. Global injectivity for odd weight.}
Suppose $[z_t]=0$ globally.  A global witness $d\in E/O$ would
satisfy $z_t(g)=(a_e(g)-1)d$ for every $g\in G_K$.
At the local element $g_0$ of Step 2 the left side is zero.
Equation (G186.6) and the invertibility of $u$ force $2d=0$.
Every global scalar $a_e(g)$ is an odd rational unit, so
$(a_e(g)-1)d=0$ for this same $d$.
Consequently $z_t$ would be identically zero globally.
Since $\xi|_{G_K}$ is nontrivial, there is a $g$ with
$\epsilon(g)=1$, forcing $t=0$.
This proves injectivity.
The original $O$-linear identification
\[
             O/2O\xrightarrow{\sim}D_e[2],
                         y\bmod2O\longmapsto y/2\bmod O
 \tag{G186.9}
\]
then proves (G186.2), without changing the coefficient field
or assuming that $O$ is unramified over $\mathbb Z_2$.

\proofstep{5. The actual ramification group and numerical lower bound.}
The host in (FG.3) contains $\mathbb Q(\zeta_7)$, so the finite
label $\eta=\xi|_{G_K}$ is an actual label, with trivial
two-power, auxiliary and Kummer-cubic factors.
Its values $\{1,-1\}$ are already in every original $O_2$.
The coefficient identity is literally
\[
 (O_2\otimes_{\mathbb Z_2}\mathbb Q_2/\mathbb Z_2)(\xi)(-e)
                         =(E_2/O_2)(\xi c_2^{-e}).
 \tag{G186.10}
\]
Both the action and the cocycle (G186.7) factor through the
actual $G_{K,\Sigma_2}$, since they are unramified outside
two and seven.  Every extra host inertia group is killed
by the cocycle and has unramified coefficient.
Thus the same local checks hold in (FC.2).
A coboundary on this quotient inflates to a coboundary on
$G_K$, which Step 4 has excluded for $t\ne0$.
Therefore (G186.2) is an injection into the actual strict group.

As a $\mathbb Z_2$-module $O$ is free of rank $b_2$.
Hence $|O/2O|=2^{b_2}$, proving (G186.3) whenever the actual
strict group is finite, as supplied by (FC.2).
For an integer cutoff $T\geq1$ there are $\lfloor T/2\rfloor$
odd positive weights below $T$.  The corresponding normalized
strict logarithms for this one fixed finite label therefore have
sum at least $\lfloor T/2\rfloor\log2$.
This is not a lower bound for their positive determinant indices:
the local comparison and degree-one torsion terms remain present.
\end{proof}

\paragraph{What the same cocycles do at even weight.}
For even positive $e$, the construction (G186.7) itself is a
global coboundary.  Write $t=y/2\bmod O$ with $y\in O$ and
take $d=-y/4\bmod O$.
Every odd rational two-adic unit $v$ satisfies $v^{-e}\equiv1$
modulo four.  If $\xi(g)=1$, then $(a_e(g)-1)d=0$; if
$\xi(g)=-1$, then $a_e(g)-1\equiv-2$ modulo four, giving
$(a_e(g)-1)d=t$.
Thus these classes vanish at even weight.
This computation makes no assertion that the entire strict
group vanishes there.

\paragraph{Scope and retained source inputs.}
The quadratic character, two-adic square, local cyclotomic element,
cocycles, full local conditions and original-ring injection are
proved directly above.  The finite Galois correspondence and
cyclotomic character convention are the standard retained ones.
The actual field, coefficient and finiteness comparison uses the
already audited (FG.3) and (FC.1)--(FC.2), with their original
local conditions; no new class computation or analytic input is used.
The finite cover $y^2=-7$ over
$\operatorname{Spec}\mathcal O_K[1/14]$ has the stated quadratic
character under the retained SGA~1, Expos\'e~V, sections 4 and 7
dictionary.  Its induced cohomology classes can have distinct
local witnesses without a global one, exactly as shown above.
The unramified-outside-two premise in AC-G182 fails for this
coefficient at seven, so its one-prime argument does not imply
vanishing here.  No theorem is asserted for general wild
$\delta$ or $\kappa$ in this note, no upper bound is inferred,
and no same-window arithmetic reduction or H04.4 conclusion
is supplied.  No Lean execution or shared-project mutation is claimed.
